\chapter{Déploiement, Tests et Sécurité} 
\minitoc
\newpage

\section{Introduction}
Tout produit logiciel professionnel, avant d'être mis entre les mains de ses utilisateurs finaux, doit passer par une phase rigoureuse de qualification et de mise en production. Ce sixième et dernier chapitre s'éloigne du développement fonctionnel pur pour aborder l'ingénierie DevOps et l'Assurance Qualité (QA). Nous détaillerons la stratégie de test mise en œuvre pour garantir l'absence de régressions, les mesures de sécurité préventives appliquées sur l'architecture backend, et enfin, le processus de déploiement (hébergement) de la base de données et de l'API Node.js dans le Cloud.

\section{Stratégie de Test et Assurance Qualité}
La complexité inhérente à une application gérant des programmes d'entraînement asynchrones et des relations inter-utilisateurs multiples rend les tests manuels insuffisants.

\subsection{Tests d'intégration avec Postman}
Pour valider le fonctionnement de notre API REST (le cœur du système), nous avons utilisé Postman de manière systématique.
Plutôt que de tester de manière ad hoc, nous avons créé une \textbf{Postman Collection} automatisée. Chaque endpoint (ex: `POST /api/programs`) a été associé à des scripts de test (écrits en JavaScript au sein de Postman) permettant de valider :
\begin{itemize}
    \item \textbf{Le code HTTP de retour :} S'assurer qu'une requête valide retourne `201 Created` ou `200 OK`, et qu'une tentative d'accès sans token JWT retourne bien `401 Unauthorized`.
    \item \textbf{Le schéma JSON :} Vérifier que l'API renvoie bien les champs attendus (ex: `program.id`, `program.days`).
    \item \textbf{Les contraintes de temps de réponse :} S'assurer que le calcul du Dashboard nutritionnel s'effectue en moins de 300 ms.
\end{itemize}

\subsection{Tests fonctionnels sur Flutter}
Côté client, le développement Flutter s'est accompagné d'une phase de tests intensifs sur divers émulateurs (Pixel 6, iPhone 14 Pro) et sur un périphérique physique.
L'objectif était de tester la résilience de l'application face à des situations dégradées :
\begin{itemize}
    \item \textbf{Perte de connexion réseau :} Vérification de l'affichage des messages d'erreur "Offline" sans crash de l'application.
    \item \textbf{Gestion des états asynchrones :} S'assurer que les indicateurs de chargement (CircularProgressIndicator) bloquent correctement les boutons pour éviter les doubles requêtes POST de la part de l'utilisateur.
\end{itemize}

\section{Sécurité de l'Application}
La manipulation de données relatives à la santé, aux mensurations et à la messagerie privée exige un niveau de sécurité optimal. Plusieurs couches de sécurité ont été implémentées.

\subsection{Sécurité de l'authentification et Hachage}
Comme abordé lors du Sprint 1, aucun mot de passe n'est stocké en clair. Le recours à \texttt{bcrypt} (avec un coût de salage à 10) rend les attaques par dictionnaire (Rainbow Tables) caduques.
La transmission de l'identifiant se fait systématiquement encapsulée dans un \textbf{Token JWT}. Ce jeton est doté d'une date d'expiration (TTL - Time To Live) courte. 

\subsection{Protection contre les attaques courantes (OWASP)}
L'API Node.js a été renforcée contre les vulnérabilités les plus fréquentes recensées par l'OWASP (Open Web Application Security Project) :
\begin{itemize}
    \item \textbf{Injections SQL :} L'utilisation exclusive de l'ORM TypeORM prévient nativement ce type d'attaque grâce à la paramétrisation automatique des requêtes. Aucune requête brute basée sur de la concaténation de chaînes utilisateur n'est exécutée.
    \item \textbf{CORS (Cross-Origin Resource Sharing) :} Un middleware restrictif a été configuré pour empêcher des domaines non autorisés de consommer l'API depuis un navigateur web. Seul le Dashboard Angular officiel est "whitelisté".
    \item \textbf{Rate Limiting :} Afin de se prémunir contre les attaques par force brute (Brute Force) sur l'endpoint de connexion, un middleware \texttt{express-rate-limit} bloque temporairement l'adresse IP source après plusieurs tentatives de login infructueuses.
\end{itemize}

\section{Processus de Déploiement (DevOps)}
Afin de rendre "CoachingApp" utilisable publiquement, les composants logiciels ont dû quitter l'environnement de développement local (`localhost`) pour être hébergés dans le Cloud.

\subsection{Hébergement de la Base de Données}
Nous avons choisi de ne pas gérer nous-mêmes le système d'exploitation du serveur de base de données. PostgreSQL a été déployé via un service géré (Managed Database as a Service) tel que Supabase ou Render.
Cela nous garantit des sauvegardes automatisées (Backups horaires), une scalabilité instantanée et un niveau de sécurité certifié, tout en nous permettant de nous concentrer sur le code métier.

\subsection{Déploiement du Serveur Node.js}
Le serveur Express a été déployé sur une instance de type \textbf{Platform as a Service (PaaS)}.
\begin{itemize}
    \item Le code source est hébergé sur un dépôt privé GitHub.
    \item \textbf{Intégration et Déploiement Continus (CI/CD) :} À chaque fois qu'un \textit{git push} est effectué vers la branche \texttt{main}, un déclencheur envoie une notification au serveur distant.
    \item Le serveur télécharge la dernière version, installe les dépendances (\texttt{npm install}), compile le TypeScript en JavaScript, puis redémarre l'instance via \texttt{PM2} (Process Manager) pour assurer zéro temps d'interruption (Zero Downtime).
\end{itemize}

\subsection{Compilation et Distribution Frontend}
Le Dashboard Web (Angular) est compilé statiquement en HTML/JS (\texttt{ng build --prod}) puis servi sur un CDN (Content Delivery Network) comme Firebase Hosting ou Vercel pour garantir des temps de chargement minimes partout dans le monde.
L'application mobile (Flutter) a été compilée en générant les \textit{Android App Bundles} (.aab) et les archives iOS (.ipa), prêtes à être publiées sur leurs magasins d'applications respectifs (Google Play Store et Apple App Store).

\section{Conclusion}
Ce huitième chapitre a démontré que notre projet n'est pas qu'une simple maquette académique, mais une solution logicielle répondant aux normes de l'industrie professionnelle. La stratégie de test assure la fiabilité des flux complexes, les mesures de sécurité garantissent la confidentialité des utilisateurs, et l'architecture de déploiement cloud permet à "CoachingApp" de supporter une charge d'utilisateurs croissante. L'ensemble de ces efforts de qualification et de mise en production vient sceller l'achèvement technique de notre projet de fin d'études.
